\lecture{8.6}{IPTW Estimation for Time-Varying Interventions}
\label{lec:8-6}

\subsection{Estimating Time-Specific Propensity Scores}

The longitudinal IPTW identification formula suggests a direct estimator. For a
fixed regime $\abar$, estimate each treatment mechanism
\[
  g_k(a_k\given \Lbar_k,\Abar_{k-1}=\abar_{k-1})
\]
and then average the weighted outcomes among individuals whose observed
treatment history matches the regime.

For two treatment times and binary treatments, estimate
\[
  g_0(a_0\given L_0)=\Prob(A_0=a_0\given L_0),
\]
and
\[
  g_1(a_1\given L_0,A_0=a_0,L_1)
  =
  \Prob(A_1=a_1\given L_0,A_0=a_0,L_1).
\]
The models for these probabilities may be logistic regressions, Super Learner
fits, or other binary regression estimators.

\subsection{Constructing Cumulative Weights}

For subject $i$, define the estimated regime probability
\[
  \gbar_{n,K,i}(\abar)
  =
  \prod_{k=0}^K
  g_{n,k}(a_k\given \Lbar_{k,i},\Abar_{k-1,i}=\abar_{k-1}).
\]
The estimated weight is
\[
  H_{n,i}(\abar)
  =
  \frac{\ind(\Abar_{K,i}=\abar_K)}{\gbar_{n,K,i}(\abar)}.
\]

\begin{estimationbox}
The IPTW estimator of the counterfactual mean under regime $\abar$ is
\[
  \psi_{n,\IPTW}(\abar)
  =
  \frac{1}{n}\sum_{i=1}^n
  H_{n,i}(\abar)Y_i.
\]
An ATE comparing regimes $\abar$ and $\abar'$ is estimated by
\[
  \psi_{n,\IPTW}(\abar)-\psi_{n,\IPTW}(\abar').
\]
\end{estimationbox}

Some analysts use stabilized weights to reduce variance. For a static regime,
a stabilized weight has a numerator that models the probability of the treatment
history using less covariate information, for example
\[
  \prod_{k=0}^K
  \Prob(A_k=a_k\given \Abar_{k-1}=\abar_{k-1}),
\]
divided by the full denominator. Stabilization changes the variance properties,
but the denominator is still the part that removes measured confounding.

\subsection{Estimating Counterfactual Means}

For $K=1$ and regime $(1,1)$,
\[
  \psi_{n,\IPTW}(1,1)
  =
  \frac{1}{n}\sum_{i=1}^n
  \frac{\ind(A_{0i}=1,A_{1i}=1)}
       {g_{n,0}(1\given L_{0i})
        g_{n,1}(1\given L_{0i},A_0=1,L_{1i})}
  Y_i.
\]
For regime $(1,0)$,
\[
  \psi_{n,\IPTW}(1,0)
  =
  \frac{1}{n}\sum_{i=1}^n
  \frac{\ind(A_{0i}=1,A_{1i}=0)}
       {g_{n,0}(1\given L_{0i})
        \{1-g_{n,1}(1\given L_{0i},A_0=1,L_{1i})\}}
  Y_i.
\]
The same pattern applies to $(0,1)$ and $(0,0)$, with the appropriate
probabilities in the denominator.

\subsection{R Implementation}

The following code implements the two-timepoint estimator in a simulated data
set with variables `A0`, `L1`, `A1`, and `Y`. The treatment at time 0 is
randomized in this example, so the first propensity model is intercept-only.

\begin{lstlisting}[language=R]
# time-0 treatment mechanism
ps0 <- glm(A0 ~ 1, family = binomial())
g_0 <- predict(ps0, type = "response")

# time-1 treatment mechanism
ps1 <- glm(A1 ~ L1 + A0, family = binomial())

# P(A1 = 1 | A0 = 1, L1 = L1_i)
g_1_A0_1 <- predict(
  ps1,
  newdata = data.frame(A0 = 1, L1 = L1),
  type = "response"
)

# P(A1 = 1 | A0 = 0, L1 = L1_i)
g_1_A0_0 <- predict(
  ps1,
  newdata = data.frame(A0 = 0, L1 = L1),
  type = "response"
)

EY_11 <- mean((A0 == 1 & A1 == 1) / (g_0 * g_1_A0_1) * Y)
EY_10 <- mean((A0 == 1 & A1 == 0) / (g_0 * (1 - g_1_A0_1)) * Y)
EY_01 <- mean((A0 == 0 & A1 == 1) / ((1 - g_0) * g_1_A0_0) * Y)
EY_00 <- mean((A0 == 0 & A1 == 0) / ((1 - g_0) * (1 - g_1_A0_0)) * Y)

ate_11_vs_00 <- EY_11 - EY_00
\end{lstlisting}

\subsection{Simulated Examples}

In the data-generating example from \cref{lec:8-5}, the true counterfactual means
are
\[
  \E\{Y(1,1)\}=6,\quad
  \E\{Y(1,0)\}=5,\quad
  \E\{Y(0,1)\}=4,\quad
  \E\{Y(0,0)\}=3.
\]
If the propensity score models are correctly specified and the sample is large,
the IPTW estimates should be close to these values. Discrepancies in finite
samples arise from random error, propensity score estimation error, and weight
instability.

\subsection{Diagnosing Extreme Weights}

Because longitudinal weights multiply inverse probabilities, diagnostics are
essential. At a minimum, inspect:
\begin{itemize}
  \item summaries of each time-specific propensity score;
  \item summaries of the cumulative denominator $\gbar_{n,K}$;
  \item summaries of the final weights $H_{n,i}(\abar)$;
  \item the number of individuals who follow each regime;
  \item whether large weights are concentrated in clinically unusual histories.
\end{itemize}

\begin{lstlisting}[language=R]
w_11 <- (A0 == 1 & A1 == 1) / (g_0 * g_1_A0_1)

summary(g_0)
summary(g_1_A0_1)
summary(w_11)
quantile(w_11, probs = c(.5, .9, .95, .99, 1))
mean(A0 == 1 & A1 == 1)
\end{lstlisting}

Large weights do not automatically imply bias, but they indicate that the
estimate depends heavily on a small number of observations. Truncating weights
can reduce variance but changes the estimand or introduces bias, so it should be
reported as a sensitivity analysis rather than hidden as a default.

\begin{keyidea}
Longitudinal IPTW is simple to compute once time-specific propensity scores are
available, but it can be highly variable because a full treatment history may be
rare even when each individual treatment decision is not rare.
\end{keyidea}
